\chapter{Ethical Issues and Data Handling Plan}

The main ethical considerations in this project come from collecting developer traces via an IDE plugin, and from the possibility that the tool could be misused as a form of monitoring. The technical work itself is not safety-critical, but the data involved can still be sensitive because it relates to how people work and may include source code artefacts.

If I collect refactoring episodes from human participants, I will use informed consent and make it clear what is logged (e.g., refactoring commands, checkpoints/snapshots, basic metadata like timestamps, and build/test outcomes) and what is \emph{not} logged (e.g., personal accounts, unrelated browsing, or continuous keystroke-level recordings unless genuinely required). I will follow a data minimisation approach: only collect the smallest amount of information needed to reconstruct refactoring steps and compute the evaluation metrics. Participants will be able to withdraw, and I will delete their data if requested.

Source code can be commercially sensitive, even when the research goal is about refactoring rather than the code itself. To avoid IP issues, I plan to run participant studies on toy projects and/or open-source repositories, rather than proprietary code. If open-source repositories are used, I will respect licensing requirements and avoid re-publishing large code excerpts in the dissertation. If I later want to evaluate on any non-public codebase, I would only do so with explicit written permission and with clear constraints on what is stored and shared.

Even without names, traces can sometimes be identifying (e.g., file paths, commit messages, timestamps, distinctive code structure). For any stored traces, I will remove or hash direct identifiers (names, usernames, machine paths, repository URLs) and avoid publishing raw logs. In the dissertation, I will present results in aggregated form where possible (e.g., across episodes), and I will be careful not to ``name and shame'' individual developers or projects when discussing low-quality trajectories.

A tool that evaluates a developer's ``process quality'' could be repurposed for workplace surveillance or performance management. Although that is not the intended use, it is important to acknowledge the risk. I will frame the tool as a developer support and coaching system (actionable feedback about code/process), not as a ranking system for people. In particular, I will avoid outputs that resemble employee scorecards or leaderboards, and I will focus on episode-level technical explanations rather than person-level judgements.

All collected data (traces and any derived metrics) will be stored securely with access limited to the project work. I will keep raw logs only for as long as needed to run experiments and compute results, and then retain only the derived, anonymised summaries required for analysis and write-up. If any data needs to be shared with a supervisor or included in an appendix, it will be anonymised and minimised to reduce the chance of leaking sensitive information.
